% This ensures VS Code compiles the main file even when you're editing a chapter.
% !TeX root = ../main.tex

\chapter{Related Works}
\label{cap:relworks}

This chapter reviews the relevant literature and previous research related to the topic
of this dissertation. The objective is to provide an overview of existing approaches,
methodologies, and technologies, highlighting their contributions, limitations, and
areas of application. By analyzing prior work, it is possible to identify current
challenges, research gaps, and opportunities that motivate the development of the
proposed solution. This review also establishes the theoretical and technical foundation
upon which the present work is built. Each of the 20 analyzed papers work will be
numbered so referencing their contents will be easier.

\section{Vision-based tracking and precision landing}

Gimbal mechanisms are primarily employed to decouple sensor orientation from the UAV's
flight dynamics. Work \#1 "Active Object Detection and Tracking Using Gimbal Mechanisms
for Autonomous Drone Applications" \cite{rel1} demonstrates that a 3-DoF gimbal
significantly reduces motion blur and peripheral distortion, improving object pose
estimation accuracy compared to passive mounts. This focus on visual stability is echoed
in work \#8 "Robust Image-Based Visual Servo Target Tracking of UAV with Depth Camera"
\cite{rel8}, which utilizes image-based visual
servoing (IBVS) to minimize tracking deviations caused by attitude perturbations.
In the context of autonomous landing, gimbals address field-of-view (FOV) limitations.
Work \#5 "Landing Control algorithm for gimbal-serviced UAVs based on field-of-view
constraints" \cite{rel5} proposes a visual servoing method that leverages gimbal agility
to ensure target features
remain within the pixel-level FOV during descent. Similarly, work \#7 "Autonomous UAV
Landing and Collision Avoidance System for Unknown Terrain Utilizing Depth Camera with
Actively Actuated Gimbal" \cite{rel7} uses an actively actuated gimbal to counteract the
narrow FOV of an
RGB-D camera, allowing the UAV to search for safe landing spots and detect obstacles
simultaneously.

\section{LiDAR integration and 3D mapping}

The integration of LiDAR sensors on gimbals enhances mapping and inspection capabilities.
Work \#4 "Autonomous Structural Inspection with an Octocopter UAV: Integration
of 3D LiDAR and Gimbal-Mounted Camera" \cite{rel4} combines a rotating 3D LiDAR for
environmental mapping with a gimballed camera for
high-resolution structural assessment. For cost-sensitive applications, Work \#20
"Designing a Compact Hexacopter with Gimballed Lidar
and Powerful Onboard Linux Computer" \cite{rel20} details a custom-built 3D-printed
gimbal designed
to stabilize a 2D LiDAR, ensuring it scans horizontally regardless of the UAV's tilt.
Several studies explore novel methods to reshape LiDAR FOV. Work \#15 "Affordable Robotic
Mobile Mapping System Based on Lidar with Additional Rotating Planar Reflector"
\cite{rel15} uses a custom
rotating mirror to extend a conic-FOV solid-state LiDAR into a 360-degree scanner.
Similarly, work \#11 "Autonomous 3D LiDAR Mapping: Transforming 1D
LiDAR to 3D LiDAR for Building Scanning with a
Mobile Robot" \cite{rel11} converts a 1D LiDAR into a 3D scanner using
an orthogonal mirror system, providing a low-cost alternative to expensive multi-beam
units. A broader overview of these technologies is provided in work \#13 "Point cloud
acquisition techniques by using scanning lidar for 3d modelling and mobile measurement"
\cite{rel13}, which
classifies scanners into mechanical, solid-state (MEMS/OPA), and multi-beam categories.

\section{Stability analysis and performance assessment}

The impact of gimbals on surveying accuracy is a critical research area.
Work \#2 "Gimbal Influence on the Stability of Exterior Orientation Parameters of UAV
Acquired Images" \cite{rel2} provides a quantitative analysis showing that low-budget
gimbals can improve orientation stability by up to 20 times in simulated conditions,
which is essential for photogrammetry without ground control points.
Work \#3 "The utilization of gimbal systems in unmanned aerial vehicles" \cite{rel3}
reviews various control algorithms (PID, Fuzzy Logic, Adaptive) used to maintain this
stability across military and civilian sectors.
For specialized mapping, Work \#12 "A light-weight laser scanner for UAV applications"
\cite{rel12} and work \#16 "Development and Performance Assessment of a Low-Cost
UAV-Borne LiDAR System (LasUAV)" \cite{rel16} evaluate the integration of
automotive-grade scanners (like IbeoLux) on UAVs. These studies conclude that the
GNSS-RTK receiver often has the greatest influence on system performance,
followed by the attitude measurement device.
Work \#17 "Airborne laser scanning, an introduction and overview" \cite{rel17} provides
the foundational context for these systems, detailing the pulse and phase-shift
measurement principles.

\section{Specialized applications and long-range detection}

Gimbals also enable reactive flight behaviors. Work \#18 "Autonomous Quadrotor
Terrain-Following with a Laser Rangefinder and Gimbal System" \cite{rel18} uses a gimbal
to point a rangefinder ahead of the UAV based on its speed, allowing real-time altitude
adjustments over uneven ground.
For ground robots in rugged terrain, work \#19 "Terrain-Adaptive Planning of a Mobile
Robot With a Multiaxis Gimbal System for Stable SLAM" \cite{rel19} introduces a unique
gimbal with Z-axis translation to dampen vertical shocks that rotational gimbals cannot
address. Finally, for counter-UAV applications, long-range detection is paramount.
Work \#9 "LiDAR Technology for UAV Detection: From Fundamentals and Operational Principles
to Advanced Detection and Classification Techniques" \cite{rel9} reviews ground-to-air
detection mechanisms.
Work \#14 "A scanning LiDAR for long range detection and tracking of UAVs" \cite{rel14}
and work \#10 "Scanning LiDAR for Telescope-based UAV Localization" \cite{rel10} both
optimize scanning patterns and optical architectures, such as using telescopes with laser
diodes, to achieve target localization at ranges up to one kilometer.
Work \#6 "UAV low-altitude obstacle detection based on the fusion of LiDAR and camera"
\cite{rel6} complements these by using an IMM-UKF algorithm to fuse LiDAR distance and
camera angle data for accurate dynamic obstacle tracking.

\section{Related works summary}

To facilitate the comparison of the reviewed studies, Table~\ref{table:relworks} summarizes
the main characteristics, methodologies, and contributions of the related works discussed
in this chapter. This overview highlights the similarities and differences among the
existing approaches, providing a clearer understanding of the current state of the art
and the research gap addressed by the proposed work. All works study the use of Computer
Vision in unmanned vehicles, as well as the use of some form of sensor fusion to acquire
depth data. The summary also serves to highlight how little the idea of gimballing a
\gls{tof} sensor has been covered in recent works due to the preference given to 3D
\gls{lidar}s. Table~\ref{table:relrefs} can be used for quick reference.

\begin{table}[ht]
  \centering
  \caption{ Article numbers and references }
  \label{table:relrefs}
  \begin{tabular}{|| c c ||} 
    \hline
    Number & Reference \\ [0.5ex] 
    \hline\hline
    [1] & \cite{rel1} \\ 
    \hline
    [2] & \cite{rel2} \\ 
    \hline
    [3] & \cite{rel3} \\ 
    \hline
    [4] & \cite{rel4} \\ 
    \hline
    [5] & \cite{rel5} \\ 
    \hline
    [6] & \cite{rel6} \\ 
    \hline
    [7] & \cite{rel7} \\ 
    \hline
    [8] & \cite{rel8} \\ 
    \hline
    [9] & \cite{rel9} \\ 
    \hline
    [10] & \cite{rel10} \\ 
    \hline
    [11] & \cite{rel11} \\ 
    \hline
    [12] & \cite{rel12} \\ 
    \hline
    [13] & \cite{rel13} \\ 
    \hline
    [14] & \cite{rel14} \\ 
    \hline
    [15] & \cite{rel15} \\ 
    \hline
    [16] & \cite{rel16} \\ 
    \hline
    [17] & \cite{rel17} \\ 
    \hline
    [18] & \cite{rel18} \\ 
    \hline
    [19] & \cite{rel19} \\ 
    \hline
    [20] & \cite{rel20} \\ 
    \hline
  \end{tabular}
\end{table}

% extra symbols for the reference table
\newcommand{\V}{ \checkmark }
\newcommand{\X}{ \text{\sffamily X} }

\begin{table}[ht]
  \centering
  \caption{ Direct comparison between all the articles used as related works }
  \label{table:relworks}
  \begin{tabular}{|| c c c c c ||} 
    \hline
    Number & Primary Sensors & Gimbal Type & Application & \begin{tabular}{@{}c@{}}Distance \\ sensor \\ on gimbal?\end{tabular} \\ [0.5ex] 
    \hline\hline
    [1] & Camera & 3-Axis & Active Tracking & ~ \\ 
    \hline
    [2] & Camera & 2-Axis & Photogrammetry & ~ \\ 
    \hline
    [3] & Camera & 3-Axis & Multi-sector Review & ~ \\ 
    \hline
    [4] & 3D \gls{lidar} + Camera & 3-Axis & Infrastructure inspection & ~ \\ 
    \hline
    [5] & Camera & 3-Axis & Autonomous Landing & ~ \\ 
    \hline
    [6] & \gls{lidar} + Camera & \X & Obstacle Avoidance & ~ \\ 
    \hline
    [7] & RGB-D Camera & 3-Axis & Landing/Avoidance & \V \\ 
    \hline
    [8] & Camera & \X & Target Tracking & ~ \\ 
    \hline
    [9] & \gls{lidar} & \X & Detection systems review & ~ \\ 
    \hline
    [10] & \gls{lidar} & 2-Axis & UAV Localization & ~ \\ 
    \hline
    [11] & \gls{tof} & 2-Axis & SLAM & ~ \\ 
    \hline
    [12] & 3D \gls{lidar} & \X & Geospatial Data & ~ \\ 
    \hline
    [13] & 3D \gls{lidar} & \X & Mapping review & ~ \\ 
    \hline
    [14] & 3D \gls{lidar} & \X & UAV Detection & ~ \\ 
    \hline
    [15] & 3D \gls{lidar} & 2-Axis & Mobile Mapping & ~ \\ 
    \hline
    [16] & 3D \gls{lidar} & \X & Forest Monitoring & ~ \\ 
    \hline
    [17] & 3D \gls{lidar} & \X & Topography & ~ \\ 
    \hline
    [18] & \gls{tof} & 2-Axis & Terrain Following & \V \\ 
    \hline
    [19] & \gls{lidar} + Camera & 2-Axis & SLAM & \V \\ 
    \hline
    [20] & 2D \gls{lidar} + Camera & 3-Axis & Sensor-Fusion & \V \\ 
    \hline
  \end{tabular}
\end{table}



